
\documentclass{article}
\usepackage{colortbl}
\usepackage{makecell}
\usepackage{multirow}
\usepackage{supertabular}

\begin{document}

\newcounter{utterance}

\centering \large Interaction Transcript for game `clean\_up', experiment `2\_hard\_5obj\_nl', episode 3 with mistral{-}large{-}3{-}azure.
\vspace{24pt}

{ \footnotesize  \setcounter{utterance}{1}
\setlength{\tabcolsep}{0pt}
\begin{supertabular}{c@{$\;$}|p{.15\linewidth}@{}p{.15\linewidth}p{.15\linewidth}p{.15\linewidth}p{.15\linewidth}p{.15\linewidth}}
    \# & $\;$A & \multicolumn{4}{c}{Game Master} & $\;\:$B\\
    \hline

    \theutterance \stepcounter{utterance}  
    & & \multicolumn{4}{p{0.6\linewidth}}{
        \cellcolor[rgb]{0.9,0.9,0.9}{
            \makecell[{{p{\linewidth}}}]{
                \texttt{\tiny{[P1$\langle$GM]}}
                \texttt{Ik ben uw spelleider, en u speelt een samenwerkingsspel met het volgende raster als speelbord:} \\
\\ 
\texttt{```} \\
\texttt{ 1234567} \\
\texttt{╔═╤═╤═══╗} \\
\texttt{║C│◌│◌TI║ 1} \\
\texttt{║◌├─┤W◌◌║ 2} \\
\texttt{╟─┤◌├─┬─╢ 3} \\
\texttt{║◌│◌│◌│◌║ 4} \\
\texttt{╟─┼─┘◌│◌║ 5} \\
\texttt{║◌│◌H┌┴─╢ 6} \\
\texttt{║◌│◌◌│◌◌║ 7} \\
\texttt{╚═╧══╧══╝} \\
\\ 
\texttt{```} \\
\\ 
\texttt{* De bovenrand toont x{-}coördinaten die naar rechts oplopen, en de rechterrand y{-}coördinaten die naar beneden oplopen.} \\
\texttt{* De volgende objecten zijn willekeurig op uw raster geplaatst: 'W', 'I', 'T', 'C', 'H'.} \\
\\ 
\\ 
\texttt{De andere speler ziet een variatie van het speelbord, waarbij de objecten op verschillende willekeurige locaties zijn geplaatst. U kunt het bord van de andere speler niet zien en zij kunnen het uwe niet zien.} \\
\\ 
\\ 
\texttt{**Doel:**} \\
\\ 
\texttt{Beide spelers moeten de objecten op hun respectievelijke achtergrond zo verplaatsen dat identieke objecten op dezelfde coördinaten terechtkomen. U moet communiceren met de andere speler om een gemeenschappelijke doelconfiguratie overeen te komen.} \\
\\ 
\\ 
\texttt{**Regels:**} \\
\\ 
\texttt{* In elke beurt kunt u precies één van de volgende twee commando's sturen:} \\
\texttt{ 1. `ZEG: <BERICHT>`: om een bericht (alles tot de volgende regeleinde) naar de andere speler te sturen. Ik zal het doorsturen naar uw partner.} \\
\texttt{2. `BEWEEG: <OBJECT>, (<X>, <Y>)`: om een object naar een nieuwe positie te verplaatsen, waarbij `<X>` de kolom is en `<Y>` de rij. Ik zal u laten weten of uw zet succesvol was of niet.} \\
\texttt{* Als u zich niet aan het formaat houdt, of meerdere commando's tegelijk verstuurt, moet ik u straffen.} \\
\texttt{* Als beide spelers meer dan 12 strafpunten verzamelen, verliezen jullie allebei het spel.} \\
\texttt{* Het is van vitaal belang dat u met de andere speler communiceert over uw doelstatus! De *enige* manier waarop u uw strategie aan de andere speler kunt overbrengen is door gebruik te maken van het `ZEG: <BERICHT>` commando!} \\
\\ 
\\ 
\texttt{**Objecten verplaatsen**} \\
\\ 
\texttt{* U kunt alleen objecten verplaatsen naar cellen binnen de grenzen van het raster. De doelcel moet leeg zijn, d.w.z. deze mag alleen het symbool '◌' bevatten.} \\
\texttt{* Als u probeert een object te verplaatsen naar een plek die niet leeg is, of als u probeert het buiten het raster te verplaatsen, moet ik u straffen. U krijgt een nieuwe poging.} \\
\texttt{* Voordat u een zet doet, moet u controleren of de doelplek leeg is, en geen letter, kader of lijn bevat!} \\
\\ 
\\ 
\texttt{**Einde van het spel**} \\
\\ 
\texttt{Als u denkt dat u het doel bereikt hebt om alle objecten uit te lijnen, kunt u de andere speler vragen het spel te beëindigen door `ZEG: klaar?` te sturen. Als de andere speler u vraagt om het spel te beëindigen en u antwoordt met `ZEG: klaar!`, zal het spel eindigen.} \\
\\ 
\texttt{Beide spelers winnen als het spel eindigt binnen 20 rondes, waarbij één ronde gedefinieerd wordt als twee spelers die elk een geldig commando sturen.} \\
\\ 
\texttt{**Score:**} \\
\\ 
\texttt{Hoe dichter de identieke objecten bij elkaar liggen op beide speelborden, hoe meer punten u krijgt. Strafpunten verminderen uw punten. Kunt u het record verslaan?} \\
\\ 
\\ 
\texttt{Stuur een bericht naar de andere speler om het spel te starten!} \\
            }
        }
    }
    & \\ \\

    \theutterance \stepcounter{utterance}  
    & \multicolumn{4}{p{0.6\linewidth}}{
        \cellcolor[rgb]{0.9,0.9,0.9}{
            \makecell[{{p{\linewidth}}}]{
                \texttt{\tiny{[P1$\rangle$GM]}}
                \texttt{Wat een leuk en uitdagend spel! Laten we beginnen met communiceren om een strategie af te spreken. Hier is mijn eerste bericht:} \\
\\ 
\texttt{```} \\
\texttt{ZEG: Hallo! Laten we eerst afspreken hoe we de objecten gaan verplaatsen. Ik stel voor dat we elk object naar een vaste coördinaat verplaatsen, bijvoorbeeld:} \\
\texttt{{-} C naar (1,1)} \\
\texttt{{-} W naar (4,2)} \\
\texttt{{-} I naar (6,1)} \\
\texttt{{-} T naar (5,1)} \\
\texttt{{-} H naar (4,6)} \\
\texttt{Wat denk jij? Of heb je een beter idee?} \\
\texttt{```} \\
            }
        }
    }
    & & \\ \\

    \theutterance \stepcounter{utterance}  
    & & & \multicolumn{2}{p{0.3\linewidth}}{
        \cellcolor[rgb]{0.9,0.9,0.9}{
            \makecell[{{p{\linewidth}}}]{
                \texttt{\tiny{[GM$|$GM]}}
                \texttt{Wat een leuk en uitdagend spel! Laten we beginnen met communiceren om een strategie af te spreken. Hier is mijn eerste bericht:} \\
\\ 
\texttt{```} \\
\texttt{ZEG: Hallo! Laten we eerst afspreken hoe we de objecten gaan verplaatsen. Ik stel voor dat we elk object naar een vaste coördinaat verplaatsen, bijvoorbeeld:} \\
\texttt{{-} C naar (1,1)} \\
\texttt{{-} W naar (4,2)} \\
\texttt{{-} I naar (6,1)} \\
\texttt{{-} T naar (5,1)} \\
\texttt{{-} H naar (4,6)} \\
\texttt{Wat denk jij? Of heb je een beter idee?} \\
\texttt{```} \\
            }
        }
    }
    & & \\ \\

    \theutterance \stepcounter{utterance}  
    & & & \multicolumn{2}{p{0.3\linewidth}}{
        \cellcolor[rgb]{0.9,0.9,0.9}{
            \makecell[{{p{\linewidth}}}]{
                \texttt{\tiny{[GM$|$GM]}}
                \texttt{Invalid format: head and tail are not empty} \\
\texttt{head: 'Wat een leuk en uitdagend spel! Laten we beginnen met communiceren om een strategie af te spreken. Hier is mijn eerste bericht:} \\
\\ 
\\ 
\texttt{'} \\
\texttt{tail: '} \\
\texttt{{-} C naar (1,1)} \\
\texttt{{-} W naar (4,2)} \\
\texttt{{-} I naar (6,1)} \\
\texttt{{-} T naar (5,1)} \\
\texttt{{-} H naar (4,6)} \\
\texttt{Wat denk jij? Of heb je een beter idee?'} \\
            }
        }
    }
    & & \\ \\

    \theutterance \stepcounter{utterance}  
    & & \multicolumn{4}{p{0.6\linewidth}}{
        \cellcolor[rgb]{0.9,0.9,0.9}{
            \makecell[{{p{\linewidth}}}]{
                \texttt{\tiny{[P1$\langle$GM]}}
                \texttt{Straf: Je bericht mag niets vóór of na het commando bevatten!} \\
\\ 
\texttt{Zorg ervoor dat je antwoord alleen `ZEG: <BERICHT>` of `BEWEEG: <OBJECT>, (<X>, <Y>)` bevat, en niets anders!} \\
\\ 
\texttt{Jullie hebben gezamenlijk 1 van 12 straffen verzameld.} \\
\texttt{Probeer het opnieuw!} \\
            }
        }
    }
    & \\ \\

    \theutterance \stepcounter{utterance}  
    & \multicolumn{4}{p{0.6\linewidth}}{
        \cellcolor[rgb]{0.9,0.9,0.9}{
            \makecell[{{p{\linewidth}}}]{
                \texttt{\tiny{[P1$\rangle$GM]}}
                \texttt{```} \\
\texttt{ZEG: Hallo! Laten we elk object naar een vaste coördinaat verplaatsen. Ik stel voor:} \\
\texttt{C (1,1), W (4,2), I (6,1), T (5,1), H (4,6).} \\
\texttt{Akkoord?} \\
\texttt{```} \\
            }
        }
    }
    & & \\ \\

    \theutterance \stepcounter{utterance}  
    & & & \multicolumn{2}{p{0.3\linewidth}}{
        \cellcolor[rgb]{0.9,0.9,0.9}{
            \makecell[{{p{\linewidth}}}]{
                \texttt{\tiny{[GM$|$GM]}}
                \texttt{```} \\
\texttt{ZEG: Hallo! Laten we elk object naar een vaste coördinaat verplaatsen. Ik stel voor:} \\
\texttt{C (1,1), W (4,2), I (6,1), T (5,1), H (4,6).} \\
\texttt{Akkoord?} \\
\texttt{```} \\
            }
        }
    }
    & & \\ \\

    \theutterance \stepcounter{utterance}  
    & & & \multicolumn{2}{p{0.3\linewidth}}{
        \cellcolor[rgb]{0.9,0.9,0.9}{
            \makecell[{{p{\linewidth}}}]{
                \texttt{\tiny{[GM$|$GM]}}
                \texttt{Invalid format: tail is not empty: '} \\
\texttt{C (1,1), W (4,2), I (6,1), T (5,1), H (4,6).} \\
\texttt{Akkoord?} \\
            }
        }
    }
    & & \\ \\

    \theutterance \stepcounter{utterance}  
    & & \multicolumn{4}{p{0.6\linewidth}}{
        \cellcolor[rgb]{0.9,0.9,0.9}{
            \makecell[{{p{\linewidth}}}]{
                \texttt{\tiny{[P1$\langle$GM]}}
                \texttt{Straf: Je bericht mag niets na het commando bevatten!} \\
\\ 
\texttt{Zorg ervoor dat je antwoord alleen `ZEG: <BERICHT>` of `BEWEEG: <OBJECT>, (<X>, <Y>)` bevat, en niets anders!} \\
\\ 
\texttt{Jullie hebben gezamenlijk 2 van 12 straffen verzameld.} \\
\texttt{Probeer het opnieuw!} \\
            }
        }
    }
    & \\ \\

    \theutterance \stepcounter{utterance}  
    & \multicolumn{4}{p{0.6\linewidth}}{
        \cellcolor[rgb]{0.9,0.9,0.9}{
            \makecell[{{p{\linewidth}}}]{
                \texttt{\tiny{[P1$\rangle$GM]}}
                \texttt{```} \\
\texttt{ZEG: Laten we vaste posities afspreken: C(1,1), W(4,2), I(6,1), T(5,1), H(4,6). Akkoord?} \\
\texttt{```} \\
            }
        }
    }
    & & \\ \\

    \theutterance \stepcounter{utterance}  
    & & & \multicolumn{2}{p{0.3\linewidth}}{
        \cellcolor[rgb]{0.9,0.9,0.9}{
            \makecell[{{p{\linewidth}}}]{
                \texttt{\tiny{[GM$|$GM]}}
                \texttt{```} \\
\texttt{ZEG: Laten we vaste posities afspreken: C(1,1), W(4,2), I(6,1), T(5,1), H(4,6). Akkoord?} \\
\texttt{```} \\
            }
        }
    }
    & & \\ \\

    \theutterance \stepcounter{utterance}  
    & & \multicolumn{4}{p{0.6\linewidth}}{
        \cellcolor[rgb]{0.9,0.9,0.9}{
            \makecell[{{p{\linewidth}}}]{
                \texttt{\tiny{[GM$\rangle$P2]}}
                \texttt{Ik ben uw spelleider, en u speelt een samenwerkingsspel met het volgende raster als speelbord:} \\
\\ 
\texttt{```} \\
\texttt{ 1234567} \\
\texttt{╔═╤═╤═══╗} \\
\texttt{║H│◌│◌◌◌║ 1} \\
\texttt{║◌├─┤◌◌◌║ 2} \\
\texttt{╟─┤◌├─┬─╢ 3} \\
\texttt{║◌│◌│C│◌║ 4} \\
\texttt{╟─┼─┘◌│◌║ 5} \\
\texttt{║◌│◌◌┌┴─╢ 6} \\
\texttt{║T│WI│◌◌║ 7} \\
\texttt{╚═╧══╧══╝} \\
\\ 
\texttt{```} \\
\\ 
\texttt{* De bovenrand toont x{-}coördinaten die naar rechts oplopen, en de rechterrand y{-}coördinaten die naar beneden oplopen.} \\
\texttt{* De volgende objecten zijn willekeurig op uw raster geplaatst: 'W', 'I', 'T', 'C', 'H'.} \\
\\ 
\\ 
\texttt{De andere speler ziet een variatie van het speelbord, waarbij de objecten op verschillende willekeurige locaties zijn geplaatst. U kunt het bord van de andere speler niet zien en zij kunnen het uwe niet zien.} \\
\\ 
\\ 
\texttt{**Doel:**} \\
\\ 
\texttt{Beide spelers moeten de objecten op hun respectievelijke achtergrond zo verplaatsen dat identieke objecten op dezelfde coördinaten terechtkomen. U moet communiceren met de andere speler om een gemeenschappelijke doelconfiguratie overeen te komen.} \\
\\ 
\\ 
\texttt{**Regels:**} \\
\\ 
\texttt{* In elke beurt kunt u precies één van de volgende twee commando's sturen:} \\
\texttt{ 1. `ZEG: <BERICHT>`: om een bericht (alles tot de volgende regeleinde) naar de andere speler te sturen. Ik zal het doorsturen naar uw partner.} \\
\texttt{2. `BEWEEG: <OBJECT>, (<X>, <Y>)`: om een object naar een nieuwe positie te verplaatsen, waarbij `<X>` de kolom is en `<Y>` de rij. Ik zal u laten weten of uw zet succesvol was of niet.} \\
\texttt{* Als u zich niet aan het formaat houdt, of meerdere commando's tegelijk verstuurt, moet ik u straffen.} \\
\texttt{* Als beide spelers meer dan 12 strafpunten verzamelen, verliezen jullie allebei het spel.} \\
\texttt{* Het is van vitaal belang dat u met de andere speler communiceert over uw doelstatus! De *enige* manier waarop u uw strategie aan de andere speler kunt overbrengen is door gebruik te maken van het `ZEG: <BERICHT>` commando!} \\
\\ 
\\ 
\texttt{**Objecten verplaatsen**} \\
\\ 
\texttt{* U kunt alleen objecten verplaatsen naar cellen binnen de grenzen van het raster. De doelcel moet leeg zijn, d.w.z. deze mag alleen het symbool '◌' bevatten.} \\
\texttt{* Als u probeert een object te verplaatsen naar een plek die niet leeg is, of als u probeert het buiten het raster te verplaatsen, moet ik u straffen. U krijgt een nieuwe poging.} \\
\texttt{* Voordat u een zet doet, moet u controleren of de doelplek leeg is, en geen letter, kader of lijn bevat!} \\
\\ 
\\ 
\texttt{**Einde van het spel**} \\
\\ 
\texttt{Als u denkt dat u het doel bereikt hebt om alle objecten uit te lijnen, kunt u de andere speler vragen het spel te beëindigen door `ZEG: klaar?` te sturen. Als de andere speler u vraagt om het spel te beëindigen en u antwoordt met `ZEG: klaar!`, zal het spel eindigen.} \\
\\ 
\texttt{Beide spelers winnen als het spel eindigt binnen 20 rondes, waarbij één ronde gedefinieerd wordt als twee spelers die elk een geldig commando sturen.} \\
\\ 
\texttt{**Score:**} \\
\\ 
\texttt{Hoe dichter de identieke objecten bij elkaar liggen op beide speelborden, hoe meer punten u krijgt. Strafpunten verminderen uw punten. Kunt u het record verslaan?} \\
\\ 
\texttt{De andere speler begon het spel door dit bericht te sturen:} \\
\\ 
\texttt{"Laten we vaste posities afspreken: C(1,1), W(4,2), I(6,1), T(5,1), H(4,6). Akkoord?"} \\
\\ 
\texttt{Wat is uw eerste opdracht?} \\
            }
        }
    }
    & \\ \\

    \theutterance \stepcounter{utterance}  
    & & & \multicolumn{4}{p{0.6\linewidth}}{
        \cellcolor[rgb]{0.9,0.9,0.9}{
            \makecell[{{p{\linewidth}}}]{
                \texttt{\tiny{[GM$\langle$P2]}}
                \texttt{Gegeven de start van de andere speler, lijkt het erop dat ze een voorstel doen voor vaste posities voor elk object. Laten we dit voorstel evalueren en bevestigen of eventueel aanpassen op basis van mijn huidige bord.} \\
\\ 
\texttt{Hier is mijn huidige bord voor referentie:} \\
\\ 
\texttt{```} \\
\texttt{ 1234567} \\
\texttt{╔═╤═╤═══╗} \\
\texttt{║H│◌│◌◌◌║ 1} \\
\texttt{║◌├─┤◌◌◌║ 2} \\
\texttt{╟─┤◌├─┬─╢ 3} \\
\texttt{║◌│◌│C│◌║ 4} \\
\texttt{╟─┼─┘◌│◌║ 5} \\
\texttt{║◌│◌◌┌┴─╢ 6} \\
\texttt{║T│WI│◌◌║ 7} \\
\texttt{╚═╧══╧══╝} \\
\texttt{```} \\
\\ 
\texttt{De voorgestelde posities zijn:} \\
\texttt{{-} C(1,1)} \\
\texttt{{-} W(4,2)} \\
\texttt{{-} I(6,1)} \\
\texttt{{-} T(5,1)} \\
\texttt{{-} H(4,6)} \\
\\ 
\texttt{Ik zal eerst kijken of deze posities op mijn bord beschikbaar (leeg) zijn:} \\
\texttt{{-} (1,1): Bevat 'H' {-}> niet leeg} \\
\texttt{{-} (4,2): Bevat '◌' {-}> leeg} \\
\texttt{{-} (6,1): Bevat '◌' {-}> leeg (maar let op: kolom 6 rij 1 lijkt leeg, maar rij 1 kolom 2 heeft een '◌')} \\
\texttt{{-} (5,1): Bevat '◌' {-}> leeg} \\
\texttt{{-} (4,6): Bevat '◌' {-}> leeg} \\
\\ 
\texttt{Ik zie dat de voorgestelde positie voor 'H' (4,6) vrij is, maar de positie voor 'C' (1,1) momenteel bezet is door 'H'. Ik zal het voorstel van de andere speler bevestigen en aangeven dat ik 'H' moet verplaatsen om ruimte te maken voor 'C'.} \\
\\ 
\texttt{Mijn eerste opdracht zal zijn om het voorstel te bevestigen en mijn eerste zet te doen om 'H' te verplaatsen naar een vrije locatie:} \\
\\ 
\texttt{```} \\
\texttt{ZEG: Akkoord met vaste posities. Ik verplaats H naar (4,6) zodat C naar (1,1) kan. Ik begin met H naar (4,6) te verplaatsen.} \\
\texttt{BEWEEG: H, (4,6)} \\
\texttt{```} \\
            }
        }
    }
     \\ \\

    \theutterance \stepcounter{utterance}  
    & & & \multicolumn{2}{p{0.3\linewidth}}{
        \cellcolor[rgb]{0.9,0.9,0.9}{
            \makecell[{{p{\linewidth}}}]{
                \texttt{\tiny{[GM$|$GM]}}
                \texttt{Gegeven de start van de andere speler, lijkt het erop dat ze een voorstel doen voor vaste posities voor elk object. Laten we dit voorstel evalueren en bevestigen of eventueel aanpassen op basis van mijn huidige bord.} \\
\\ 
\texttt{Hier is mijn huidige bord voor referentie:} \\
\\ 
\texttt{```} \\
\texttt{ 1234567} \\
\texttt{╔═╤═╤═══╗} \\
\texttt{║H│◌│◌◌◌║ 1} \\
\texttt{║◌├─┤◌◌◌║ 2} \\
\texttt{╟─┤◌├─┬─╢ 3} \\
\texttt{║◌│◌│C│◌║ 4} \\
\texttt{╟─┼─┘◌│◌║ 5} \\
\texttt{║◌│◌◌┌┴─╢ 6} \\
\texttt{║T│WI│◌◌║ 7} \\
\texttt{╚═╧══╧══╝} \\
\texttt{```} \\
\\ 
\texttt{De voorgestelde posities zijn:} \\
\texttt{{-} C(1,1)} \\
\texttt{{-} W(4,2)} \\
\texttt{{-} I(6,1)} \\
\texttt{{-} T(5,1)} \\
\texttt{{-} H(4,6)} \\
\\ 
\texttt{Ik zal eerst kijken of deze posities op mijn bord beschikbaar (leeg) zijn:} \\
\texttt{{-} (1,1): Bevat 'H' {-}> niet leeg} \\
\texttt{{-} (4,2): Bevat '◌' {-}> leeg} \\
\texttt{{-} (6,1): Bevat '◌' {-}> leeg (maar let op: kolom 6 rij 1 lijkt leeg, maar rij 1 kolom 2 heeft een '◌')} \\
\texttt{{-} (5,1): Bevat '◌' {-}> leeg} \\
\texttt{{-} (4,6): Bevat '◌' {-}> leeg} \\
\\ 
\texttt{Ik zie dat de voorgestelde positie voor 'H' (4,6) vrij is, maar de positie voor 'C' (1,1) momenteel bezet is door 'H'. Ik zal het voorstel van de andere speler bevestigen en aangeven dat ik 'H' moet verplaatsen om ruimte te maken voor 'C'.} \\
\\ 
\texttt{Mijn eerste opdracht zal zijn om het voorstel te bevestigen en mijn eerste zet te doen om 'H' te verplaatsen naar een vrije locatie:} \\
\\ 
\texttt{```} \\
\texttt{ZEG: Akkoord met vaste posities. Ik verplaats H naar (4,6) zodat C naar (1,1) kan. Ik begin met H naar (4,6) te verplaatsen.} \\
\texttt{BEWEEG: H, (4,6)} \\
\texttt{```} \\
            }
        }
    }
    & & \\ \\

    \theutterance \stepcounter{utterance}  
    & & & \multicolumn{2}{p{0.3\linewidth}}{
        \cellcolor[rgb]{0.9,0.9,0.9}{
            \makecell[{{p{\linewidth}}}]{
                \texttt{\tiny{[GM$|$GM]}}
                \texttt{Invalid response format: Gegeven de start van de andere speler, lijkt het erop dat ze een voorstel doen voor vaste posities voor elk object. Laten we dit voorstel evalueren en bevestigen of eventueel aanpassen op basis van mijn huidige bord.} \\
\\ 
\texttt{Hier is mijn huidige bord voor referentie:} \\
\\ 
\\ 
\texttt{ 1234567} \\
\texttt{╔═╤═╤═══╗} \\
\texttt{║H│◌│◌◌◌║ 1} \\
\texttt{║◌├─┤◌◌◌║ 2} \\
\texttt{╟─┤◌├─┬─╢ 3} \\
\texttt{║◌│◌│C│◌║ 4} \\
\texttt{╟─┼─┘◌│◌║ 5} \\
\texttt{║◌│◌◌┌┴─╢ 6} \\
\texttt{║T│WI│◌◌║ 7} \\
\texttt{╚═╧══╧══╝} \\
\\ 
\\ 
\texttt{De voorgestelde posities zijn:} \\
\texttt{{-} C(1,1)} \\
\texttt{{-} W(4,2)} \\
\texttt{{-} I(6,1)} \\
\texttt{{-} T(5,1)} \\
\texttt{{-} H(4,6)} \\
\\ 
\texttt{Ik zal eerst kijken of deze posities op mijn bord beschikbaar (leeg) zijn:} \\
\texttt{{-} (1,1): Bevat 'H' {-}> niet leeg} \\
\texttt{{-} (4,2): Bevat '◌' {-}> leeg} \\
\texttt{{-} (6,1): Bevat '◌' {-}> leeg (maar let op: kolom 6 rij 1 lijkt leeg, maar rij 1 kolom 2 heeft een '◌')} \\
\texttt{{-} (5,1): Bevat '◌' {-}> leeg} \\
\texttt{{-} (4,6): Bevat '◌' {-}> leeg} \\
\\ 
\texttt{Ik zie dat de voorgestelde positie voor 'H' (4,6) vrij is, maar de positie voor 'C' (1,1) momenteel bezet is door 'H'. Ik zal het voorstel van de andere speler bevestigen en aangeven dat ik 'H' moet verplaatsen om ruimte te maken voor 'C'.} \\
\\ 
\texttt{Mijn eerste opdracht zal zijn om het voorstel te bevestigen en mijn eerste zet te doen om 'H' te verplaatsen naar een vrije locatie:} \\
\\ 
\\ 
\texttt{ZEG: Akkoord met vaste posities. Ik verplaats H naar (4,6) zodat C naar (1,1) kan. Ik begin met H naar (4,6) te verplaatsen.} \\
\texttt{BEWEEG: H, (4,6)} \\
            }
        }
    }
    & & \\ \\

    \theutterance \stepcounter{utterance}  
    & & \multicolumn{4}{p{0.6\linewidth}}{
        \cellcolor[rgb]{0.9,0.9,0.9}{
            \makecell[{{p{\linewidth}}}]{
                \texttt{\tiny{[GM$\rangle$P2]}}
                \texttt{Straf: Je bericht bevat meer dan één commando!} \\
\\ 
\texttt{Zorg ervoor dat je antwoord alleen `ZEG: <BERICHT>` of `BEWEEG: <OBJECT>, (<X>, <Y>)` bevat, en niets anders!} \\
\\ 
\texttt{Jullie hebben gezamenlijk 3 van 12 straffen verzameld.} \\
\texttt{Probeer het opnieuw!} \\
            }
        }
    }
    & \\ \\

    \theutterance \stepcounter{utterance}  
    & & & \multicolumn{4}{p{0.6\linewidth}}{
        \cellcolor[rgb]{0.9,0.9,0.9}{
            \makecell[{{p{\linewidth}}}]{
                \texttt{\tiny{[GM$\langle$P2]}}
                \texttt{Ik begrijp het, ik zal één commando per beurt sturen.} \\
\\ 
\texttt{```} \\
\texttt{ZEG: Akkoord met vaste posities. Ik begin met H naar (4,6) verplaatsen.} \\
\texttt{```} \\
            }
        }
    }
     \\ \\

    \theutterance \stepcounter{utterance}  
    & & & \multicolumn{2}{p{0.3\linewidth}}{
        \cellcolor[rgb]{0.9,0.9,0.9}{
            \makecell[{{p{\linewidth}}}]{
                \texttt{\tiny{[GM$|$GM]}}
                \texttt{Ik begrijp het, ik zal één commando per beurt sturen.} \\
\\ 
\texttt{```} \\
\texttt{ZEG: Akkoord met vaste posities. Ik begin met H naar (4,6) verplaatsen.} \\
\texttt{```} \\
            }
        }
    }
    & & \\ \\

    \theutterance \stepcounter{utterance}  
    & & & \multicolumn{2}{p{0.3\linewidth}}{
        \cellcolor[rgb]{0.9,0.9,0.9}{
            \makecell[{{p{\linewidth}}}]{
                \texttt{\tiny{[GM$|$GM]}}
                \texttt{Invalid format: head is not empty: 'Ik begrijp het, ik zal één commando per beurt sturen.} \\
\\ 
\\ 
\texttt{'} \\
            }
        }
    }
    & & \\ \\

    \theutterance \stepcounter{utterance}  
    & & \multicolumn{4}{p{0.6\linewidth}}{
        \cellcolor[rgb]{0.9,0.9,0.9}{
            \makecell[{{p{\linewidth}}}]{
                \texttt{\tiny{[GM$\rangle$P2]}}
                \texttt{Straf: Je bericht mag niets vóór het commando bevatten!} \\
\\ 
\texttt{Zorg ervoor dat je antwoord alleen `ZEG: <BERICHT>` of `BEWEEG: <OBJECT>, (<X>, <Y>)` bevat, en niets anders!} \\
\\ 
\texttt{Jullie hebben gezamenlijk 4 van 12 straffen verzameld.} \\
\texttt{Probeer het opnieuw!} \\
            }
        }
    }
    & \\ \\

    \theutterance \stepcounter{utterance}  
    & & & \multicolumn{4}{p{0.6\linewidth}}{
        \cellcolor[rgb]{0.9,0.9,0.9}{
            \makecell[{{p{\linewidth}}}]{
                \texttt{\tiny{[GM$\langle$P2]}}
                \texttt{```} \\
\texttt{ZEG: Akkoord met C(1,1), W(4,2), I(6,1), T(5,1), H(4,6). Ik verplaats H eerst naar (4,6).} \\
\texttt{```} \\
            }
        }
    }
     \\ \\

    \theutterance \stepcounter{utterance}  
    & & & \multicolumn{2}{p{0.3\linewidth}}{
        \cellcolor[rgb]{0.9,0.9,0.9}{
            \makecell[{{p{\linewidth}}}]{
                \texttt{\tiny{[GM$|$GM]}}
                \texttt{```} \\
\texttt{ZEG: Akkoord met C(1,1), W(4,2), I(6,1), T(5,1), H(4,6). Ik verplaats H eerst naar (4,6).} \\
\texttt{```} \\
            }
        }
    }
    & & \\ \\

    \theutterance \stepcounter{utterance}  
    & & \multicolumn{4}{p{0.6\linewidth}}{
        \cellcolor[rgb]{0.9,0.9,0.9}{
            \makecell[{{p{\linewidth}}}]{
                \texttt{\tiny{[P1$\langle$GM]}}
                \texttt{Je bericht is doorgestuurd naar de andere speler.} \\
\\ 
\\ 
\texttt{Je speelt momenteel ronde 1 van maximaal 20.} \\
\texttt{Jullie hebben gezamenlijk 4 van 12 straffen verzameld.} \\
\texttt{De andere speler stuurde dit bericht:} \\
\\ 
\\ 
\texttt{"Akkoord met C(1,1), W(4,2), I(6,1), T(5,1), H(4,6). Ik verplaats H eerst naar (4,6)."} \\
\\ 
\\ 
\texttt{Wat is je volgende opdracht?} \\
            }
        }
    }
    & \\ \\

    \theutterance \stepcounter{utterance}  
    & \multicolumn{4}{p{0.6\linewidth}}{
        \cellcolor[rgb]{0.9,0.9,0.9}{
            \makecell[{{p{\linewidth}}}]{
                \texttt{\tiny{[P1$\rangle$GM]}}
                \texttt{```} \\
\texttt{BEWEEG: C, (1, 1)} \\
\texttt{```} \\
            }
        }
    }
    & & \\ \\

    \theutterance \stepcounter{utterance}  
    & & & \multicolumn{2}{p{0.3\linewidth}}{
        \cellcolor[rgb]{0.9,0.9,0.9}{
            \makecell[{{p{\linewidth}}}]{
                \texttt{\tiny{[GM$|$GM]}}
                \texttt{```} \\
\texttt{BEWEEG: C, (1, 1)} \\
\texttt{```} \\
            }
        }
    }
    & & \\ \\

    \theutterance \stepcounter{utterance}  
    & & & \multicolumn{2}{p{0.3\linewidth}}{
        \cellcolor[rgb]{0.9,0.9,0.9}{
            \makecell[{{p{\linewidth}}}]{
                \texttt{\tiny{[GM$|$GM]}}
                \texttt{Straf: (1,1) is niet leeg, maar bevat 'C'.} \\
\\ 
\texttt{Jullie hebben gezamenlijk 5 van 12 straffen verzameld.} \\
\texttt{Probeer het opnieuw!} \\
            }
        }
    }
    & & \\ \\

    \theutterance \stepcounter{utterance}  
    & & \multicolumn{4}{p{0.6\linewidth}}{
        \cellcolor[rgb]{0.9,0.9,0.9}{
            \makecell[{{p{\linewidth}}}]{
                \texttt{\tiny{[P1$\langle$GM]}}
                \texttt{Straf: (1,1) is niet leeg, maar bevat 'C'.} \\
\\ 
\texttt{Jullie hebben gezamenlijk 5 van 12 straffen verzameld.} \\
\texttt{Probeer het opnieuw!} \\
            }
        }
    }
    & \\ \\

    \theutterance \stepcounter{utterance}  
    & \multicolumn{4}{p{0.6\linewidth}}{
        \cellcolor[rgb]{0.9,0.9,0.9}{
            \makecell[{{p{\linewidth}}}]{
                \texttt{\tiny{[P1$\rangle$GM]}}
                \texttt{```} \\
\texttt{BEWEEG: W, (4, 2)} \\
\texttt{```} \\
            }
        }
    }
    & & \\ \\

    \theutterance \stepcounter{utterance}  
    & & & \multicolumn{2}{p{0.3\linewidth}}{
        \cellcolor[rgb]{0.9,0.9,0.9}{
            \makecell[{{p{\linewidth}}}]{
                \texttt{\tiny{[GM$|$GM]}}
                \texttt{```} \\
\texttt{BEWEEG: W, (4, 2)} \\
\texttt{```} \\
            }
        }
    }
    & & \\ \\

    \theutterance \stepcounter{utterance}  
    & & & \multicolumn{2}{p{0.3\linewidth}}{
        \cellcolor[rgb]{0.9,0.9,0.9}{
            \makecell[{{p{\linewidth}}}]{
                \texttt{\tiny{[GM$|$GM]}}
                \texttt{Straf: (4,2) is niet leeg, maar bevat '┤'.} \\
\\ 
\texttt{Jullie hebben gezamenlijk 6 van 12 straffen verzameld.} \\
\texttt{Probeer het opnieuw!} \\
            }
        }
    }
    & & \\ \\

    \theutterance \stepcounter{utterance}  
    & & \multicolumn{4}{p{0.6\linewidth}}{
        \cellcolor[rgb]{0.9,0.9,0.9}{
            \makecell[{{p{\linewidth}}}]{
                \texttt{\tiny{[P1$\langle$GM]}}
                \texttt{Straf: (4,2) is niet leeg, maar bevat '┤'.} \\
\\ 
\texttt{Jullie hebben gezamenlijk 6 van 12 straffen verzameld.} \\
\texttt{Probeer het opnieuw!} \\
            }
        }
    }
    & \\ \\

    \theutterance \stepcounter{utterance}  
    & \multicolumn{4}{p{0.6\linewidth}}{
        \cellcolor[rgb]{0.9,0.9,0.9}{
            \makecell[{{p{\linewidth}}}]{
                \texttt{\tiny{[P1$\rangle$GM]}}
                \texttt{```} \\
\texttt{BEWEEG: I, (6, 1)} \\
\texttt{```} \\
            }
        }
    }
    & & \\ \\

    \theutterance \stepcounter{utterance}  
    & & & \multicolumn{2}{p{0.3\linewidth}}{
        \cellcolor[rgb]{0.9,0.9,0.9}{
            \makecell[{{p{\linewidth}}}]{
                \texttt{\tiny{[GM$|$GM]}}
                \texttt{```} \\
\texttt{BEWEEG: I, (6, 1)} \\
\texttt{```} \\
            }
        }
    }
    & & \\ \\

    \theutterance \stepcounter{utterance}  
    & & & \multicolumn{2}{p{0.3\linewidth}}{
        \cellcolor[rgb]{0.9,0.9,0.9}{
            \makecell[{{p{\linewidth}}}]{
                \texttt{\tiny{[GM$|$GM]}}
                \texttt{Straf: (6,1) is niet leeg, maar bevat 'T'.} \\
\\ 
\texttt{Jullie hebben gezamenlijk 7 van 12 straffen verzameld.} \\
\texttt{Probeer het opnieuw!} \\
            }
        }
    }
    & & \\ \\

    \theutterance \stepcounter{utterance}  
    & & \multicolumn{4}{p{0.6\linewidth}}{
        \cellcolor[rgb]{0.9,0.9,0.9}{
            \makecell[{{p{\linewidth}}}]{
                \texttt{\tiny{[P1$\langle$GM]}}
                \texttt{Straf: (6,1) is niet leeg, maar bevat 'T'.} \\
\\ 
\texttt{Jullie hebben gezamenlijk 7 van 12 straffen verzameld.} \\
\texttt{Probeer het opnieuw!} \\
            }
        }
    }
    & \\ \\

    \theutterance \stepcounter{utterance}  
    & \multicolumn{4}{p{0.6\linewidth}}{
        \cellcolor[rgb]{0.9,0.9,0.9}{
            \makecell[{{p{\linewidth}}}]{
                \texttt{\tiny{[P1$\rangle$GM]}}
                \texttt{```} \\
\texttt{ZEG: Mijn bord heeft lijnen op (4,2) en (6,1). Laten we alternatieve posities kiezen:} \\
\texttt{C(1,2), W(3,2), I(6,3), T(5,3), H(4,5). Akkoord?} \\
\texttt{```} \\
            }
        }
    }
    & & \\ \\

    \theutterance \stepcounter{utterance}  
    & & & \multicolumn{2}{p{0.3\linewidth}}{
        \cellcolor[rgb]{0.9,0.9,0.9}{
            \makecell[{{p{\linewidth}}}]{
                \texttt{\tiny{[GM$|$GM]}}
                \texttt{```} \\
\texttt{ZEG: Mijn bord heeft lijnen op (4,2) en (6,1). Laten we alternatieve posities kiezen:} \\
\texttt{C(1,2), W(3,2), I(6,3), T(5,3), H(4,5). Akkoord?} \\
\texttt{```} \\
            }
        }
    }
    & & \\ \\

    \theutterance \stepcounter{utterance}  
    & & & \multicolumn{2}{p{0.3\linewidth}}{
        \cellcolor[rgb]{0.9,0.9,0.9}{
            \makecell[{{p{\linewidth}}}]{
                \texttt{\tiny{[GM$|$GM]}}
                \texttt{Invalid format: tail is not empty: '} \\
\texttt{C(1,2), W(3,2), I(6,3), T(5,3), H(4,5). Akkoord?} \\
            }
        }
    }
    & & \\ \\

    \theutterance \stepcounter{utterance}  
    & & \multicolumn{4}{p{0.6\linewidth}}{
        \cellcolor[rgb]{0.9,0.9,0.9}{
            \makecell[{{p{\linewidth}}}]{
                \texttt{\tiny{[P1$\langle$GM]}}
                \texttt{Straf: Je bericht mag niets na het commando bevatten!} \\
\\ 
\texttt{Zorg ervoor dat je antwoord alleen `ZEG: <BERICHT>` of `BEWEEG: <OBJECT>, (<X>, <Y>)` bevat, en niets anders!} \\
\\ 
\texttt{Jullie hebben gezamenlijk 8 van 12 straffen verzameld.} \\
\texttt{Probeer het opnieuw!} \\
            }
        }
    }
    & \\ \\

    \theutterance \stepcounter{utterance}  
    & \multicolumn{4}{p{0.6\linewidth}}{
        \cellcolor[rgb]{0.9,0.9,0.9}{
            \makecell[{{p{\linewidth}}}]{
                \texttt{\tiny{[P1$\rangle$GM]}}
                \texttt{```} \\
\texttt{ZEG: Mijn bord blokkeert (4,2) en (6,1). Nieuwe voorstellen:} \\
\texttt{C(1,2), W(3,2), I(6,3), T(5,3), H(4,5). Akkoord?} \\
\texttt{```} \\
            }
        }
    }
    & & \\ \\

    \theutterance \stepcounter{utterance}  
    & & & \multicolumn{2}{p{0.3\linewidth}}{
        \cellcolor[rgb]{0.9,0.9,0.9}{
            \makecell[{{p{\linewidth}}}]{
                \texttt{\tiny{[GM$|$GM]}}
                \texttt{```} \\
\texttt{ZEG: Mijn bord blokkeert (4,2) en (6,1). Nieuwe voorstellen:} \\
\texttt{C(1,2), W(3,2), I(6,3), T(5,3), H(4,5). Akkoord?} \\
\texttt{```} \\
            }
        }
    }
    & & \\ \\

    \theutterance \stepcounter{utterance}  
    & & & \multicolumn{2}{p{0.3\linewidth}}{
        \cellcolor[rgb]{0.9,0.9,0.9}{
            \makecell[{{p{\linewidth}}}]{
                \texttt{\tiny{[GM$|$GM]}}
                \texttt{Invalid format: tail is not empty: '} \\
\texttt{C(1,2), W(3,2), I(6,3), T(5,3), H(4,5). Akkoord?} \\
            }
        }
    }
    & & \\ \\

    \theutterance \stepcounter{utterance}  
    & & \multicolumn{4}{p{0.6\linewidth}}{
        \cellcolor[rgb]{0.9,0.9,0.9}{
            \makecell[{{p{\linewidth}}}]{
                \texttt{\tiny{[P1$\langle$GM]}}
                \texttt{Straf: Je bericht mag niets na het commando bevatten!} \\
\\ 
\texttt{Zorg ervoor dat je antwoord alleen `ZEG: <BERICHT>` of `BEWEEG: <OBJECT>, (<X>, <Y>)` bevat, en niets anders!} \\
\\ 
\texttt{Jullie hebben gezamenlijk 9 van 12 straffen verzameld.} \\
\texttt{Probeer het opnieuw!} \\
            }
        }
    }
    & \\ \\

    \theutterance \stepcounter{utterance}  
    & \multicolumn{4}{p{0.6\linewidth}}{
        \cellcolor[rgb]{0.9,0.9,0.9}{
            \makecell[{{p{\linewidth}}}]{
                \texttt{\tiny{[P1$\rangle$GM]}}
                \texttt{```} \\
\texttt{ZEG: Blokkades op (4,2) en (6,1). Alternatieven:} \\
\texttt{C(1,2) W(3,2) I(6,3) T(5,3) H(4,5)} \\
\texttt{```} \\
            }
        }
    }
    & & \\ \\

    \theutterance \stepcounter{utterance}  
    & & & \multicolumn{2}{p{0.3\linewidth}}{
        \cellcolor[rgb]{0.9,0.9,0.9}{
            \makecell[{{p{\linewidth}}}]{
                \texttt{\tiny{[GM$|$GM]}}
                \texttt{```} \\
\texttt{ZEG: Blokkades op (4,2) en (6,1). Alternatieven:} \\
\texttt{C(1,2) W(3,2) I(6,3) T(5,3) H(4,5)} \\
\texttt{```} \\
            }
        }
    }
    & & \\ \\

    \theutterance \stepcounter{utterance}  
    & & & \multicolumn{2}{p{0.3\linewidth}}{
        \cellcolor[rgb]{0.9,0.9,0.9}{
            \makecell[{{p{\linewidth}}}]{
                \texttt{\tiny{[GM$|$GM]}}
                \texttt{Invalid format: tail is not empty: '} \\
\texttt{C(1,2) W(3,2) I(6,3) T(5,3) H(4,5)} \\
            }
        }
    }
    & & \\ \\

    \theutterance \stepcounter{utterance}  
    & & \multicolumn{4}{p{0.6\linewidth}}{
        \cellcolor[rgb]{0.9,0.9,0.9}{
            \makecell[{{p{\linewidth}}}]{
                \texttt{\tiny{[P1$\langle$GM]}}
                \texttt{Straf: Je bericht mag niets na het commando bevatten!} \\
\\ 
\texttt{Zorg ervoor dat je antwoord alleen `ZEG: <BERICHT>` of `BEWEEG: <OBJECT>, (<X>, <Y>)` bevat, en niets anders!} \\
\\ 
\texttt{Jullie hebben gezamenlijk 10 van 12 straffen verzameld.} \\
\texttt{Probeer het opnieuw!} \\
            }
        }
    }
    & \\ \\

    \theutterance \stepcounter{utterance}  
    & \multicolumn{4}{p{0.6\linewidth}}{
        \cellcolor[rgb]{0.9,0.9,0.9}{
            \makecell[{{p{\linewidth}}}]{
                \texttt{\tiny{[P1$\rangle$GM]}}
                \texttt{```} \\
\texttt{ZEG: Laten we deze posities proberen:} \\
\texttt{C(1,2)} \\
\texttt{W(3,2)} \\
\texttt{I(6,3)} \\
\texttt{T(5,3)} \\
\texttt{H(4,5)} \\
\texttt{```} \\
            }
        }
    }
    & & \\ \\

    \theutterance \stepcounter{utterance}  
    & & & \multicolumn{2}{p{0.3\linewidth}}{
        \cellcolor[rgb]{0.9,0.9,0.9}{
            \makecell[{{p{\linewidth}}}]{
                \texttt{\tiny{[GM$|$GM]}}
                \texttt{```} \\
\texttt{ZEG: Laten we deze posities proberen:} \\
\texttt{C(1,2)} \\
\texttt{W(3,2)} \\
\texttt{I(6,3)} \\
\texttt{T(5,3)} \\
\texttt{H(4,5)} \\
\texttt{```} \\
            }
        }
    }
    & & \\ \\

    \theutterance \stepcounter{utterance}  
    & & & \multicolumn{2}{p{0.3\linewidth}}{
        \cellcolor[rgb]{0.9,0.9,0.9}{
            \makecell[{{p{\linewidth}}}]{
                \texttt{\tiny{[GM$|$GM]}}
                \texttt{Invalid format: tail is not empty: '} \\
\texttt{C(1,2)} \\
\texttt{W(3,2)} \\
\texttt{I(6,3)} \\
\texttt{T(5,3)} \\
\texttt{H(4,5)} \\
            }
        }
    }
    & & \\ \\

    \theutterance \stepcounter{utterance}  
    & & \multicolumn{4}{p{0.6\linewidth}}{
        \cellcolor[rgb]{0.9,0.9,0.9}{
            \makecell[{{p{\linewidth}}}]{
                \texttt{\tiny{[P1$\langle$GM]}}
                \texttt{Straf: Je bericht mag niets na het commando bevatten!} \\
\\ 
\texttt{Zorg ervoor dat je antwoord alleen `ZEG: <BERICHT>` of `BEWEEG: <OBJECT>, (<X>, <Y>)` bevat, en niets anders!} \\
\\ 
\texttt{Jullie hebben gezamenlijk 11 van 12 straffen verzameld.} \\
\texttt{Probeer het opnieuw!} \\
            }
        }
    }
    & \\ \\

    \theutterance \stepcounter{utterance}  
    & \multicolumn{4}{p{0.6\linewidth}}{
        \cellcolor[rgb]{0.9,0.9,0.9}{
            \makecell[{{p{\linewidth}}}]{
                \texttt{\tiny{[P1$\rangle$GM]}}
                \texttt{```} \\
\texttt{ZEG: C(1,2) W(3,2) I(6,3) T(5,3) H(4,5)} \\
\texttt{```} \\
            }
        }
    }
    & & \\ \\

    \theutterance \stepcounter{utterance}  
    & & & \multicolumn{2}{p{0.3\linewidth}}{
        \cellcolor[rgb]{0.9,0.9,0.9}{
            \makecell[{{p{\linewidth}}}]{
                \texttt{\tiny{[GM$|$GM]}}
                \texttt{```} \\
\texttt{ZEG: C(1,2) W(3,2) I(6,3) T(5,3) H(4,5)} \\
\texttt{```} \\
            }
        }
    }
    & & \\ \\

    \theutterance \stepcounter{utterance}  
    & & \multicolumn{4}{p{0.6\linewidth}}{
        \cellcolor[rgb]{0.9,0.9,0.9}{
            \makecell[{{p{\linewidth}}}]{
                \texttt{\tiny{[GM$\rangle$P2]}}
                \texttt{Je bericht is doorgestuurd naar de andere speler.} \\
\\ 
\\ 
\texttt{Je speelt momenteel ronde 2 van maximaal 20.} \\
\texttt{Jullie hebben gezamenlijk 11 van 12 straffen verzameld.} \\
\texttt{De andere speler stuurde dit bericht:} \\
\\ 
\\ 
\texttt{"C(1,2) W(3,2) I(6,3) T(5,3) H(4,5)"} \\
\\ 
\\ 
\texttt{Wat is je volgende opdracht?} \\
            }
        }
    }
    & \\ \\

    \theutterance \stepcounter{utterance}  
    & & & \multicolumn{4}{p{0.6\linewidth}}{
        \cellcolor[rgb]{0.9,0.9,0.9}{
            \makecell[{{p{\linewidth}}}]{
                \texttt{\tiny{[GM$\langle$P2]}}
                \texttt{De andere speler heeft een nieuw voorstel gedaan:} \\
\texttt{{-} **C(1,2)**} \\
\texttt{{-} **W(3,2)**} \\
\texttt{{-} **I(6,3)**} \\
\texttt{{-} **T(5,3)**} \\
\texttt{{-} **H(4,5)**} \\
\\ 
\texttt{Ik controleer eerst of deze posities op mijn bord leeg zijn:} \\
\texttt{{-} **(1,2)**: Bevat **◌** (leeg) → OK} \\
\texttt{{-} **(3,2)**: Bevat **◌** (leeg) → OK} \\
\texttt{{-} **(6,3)**: Bevat **◌** (leeg) → OK} \\
\texttt{{-} **(5,3)**: Bevat **◌** (leeg) → OK} \\
\texttt{{-} **(4,5)**: Bevat **◌** (leeg) → OK} \\
\\ 
\texttt{Dit voorstel lijkt haalbaar. Ik zal eerst **C** verplaatsen naar **(1,2)**, omdat dat een logische eerste stap is.} \\
\\ 
\texttt{```} \\
\texttt{BEWEEG: C, (1,2)} \\
\texttt{```} \\
            }
        }
    }
     \\ \\

    \theutterance \stepcounter{utterance}  
    & & & \multicolumn{2}{p{0.3\linewidth}}{
        \cellcolor[rgb]{0.9,0.9,0.9}{
            \makecell[{{p{\linewidth}}}]{
                \texttt{\tiny{[GM$|$GM]}}
                \texttt{De andere speler heeft een nieuw voorstel gedaan:} \\
\texttt{{-} **C(1,2)**} \\
\texttt{{-} **W(3,2)**} \\
\texttt{{-} **I(6,3)**} \\
\texttt{{-} **T(5,3)**} \\
\texttt{{-} **H(4,5)**} \\
\\ 
\texttt{Ik controleer eerst of deze posities op mijn bord leeg zijn:} \\
\texttt{{-} **(1,2)**: Bevat **◌** (leeg) → OK} \\
\texttt{{-} **(3,2)**: Bevat **◌** (leeg) → OK} \\
\texttt{{-} **(6,3)**: Bevat **◌** (leeg) → OK} \\
\texttt{{-} **(5,3)**: Bevat **◌** (leeg) → OK} \\
\texttt{{-} **(4,5)**: Bevat **◌** (leeg) → OK} \\
\\ 
\texttt{Dit voorstel lijkt haalbaar. Ik zal eerst **C** verplaatsen naar **(1,2)**, omdat dat een logische eerste stap is.} \\
\\ 
\texttt{```} \\
\texttt{BEWEEG: C, (1,2)} \\
\texttt{```} \\
            }
        }
    }
    & & \\ \\

    \theutterance \stepcounter{utterance}  
    & & & \multicolumn{2}{p{0.3\linewidth}}{
        \cellcolor[rgb]{0.9,0.9,0.9}{
            \makecell[{{p{\linewidth}}}]{
                \texttt{\tiny{[GM$|$GM]}}
                \texttt{Invalid format: head is not empty: 'De andere speler heeft een nieuw voorstel gedaan:} \\
\texttt{{-} **C(1,2)**} \\
\texttt{{-} **W(3,2)**} \\
\texttt{{-} **I(6,3)**} \\
\texttt{{-} **T(5,3)**} \\
\texttt{{-} **H(4,5)**} \\
\\ 
\texttt{Ik controleer eerst of deze posities op mijn bord leeg zijn:} \\
\texttt{{-} **(1,2)**: Bevat **◌** (leeg) → OK} \\
\texttt{{-} **(3,2)**: Bevat **◌** (leeg) → OK} \\
\texttt{{-} **(6,3)**: Bevat **◌** (leeg) → OK} \\
\texttt{{-} **(5,3)**: Bevat **◌** (leeg) → OK} \\
\texttt{{-} **(4,5)**: Bevat **◌** (leeg) → OK} \\
\\ 
\texttt{Dit voorstel lijkt haalbaar. Ik zal eerst **C** verplaatsen naar **(1,2)**, omdat dat een logische eerste stap is.} \\
\\ 
\\ 
\texttt{'} \\
            }
        }
    }
    & & \\ \\

    \theutterance \stepcounter{utterance}  
    & & \multicolumn{4}{p{0.6\linewidth}}{
        \cellcolor[rgb]{0.9,0.9,0.9}{
            \makecell[{{p{\linewidth}}}]{
                \texttt{\tiny{[GM$\rangle$P2]}}
                \texttt{Straf: Je bericht mag niets vóór het commando bevatten!} \\
\\ 
\texttt{Zorg ervoor dat je antwoord alleen `ZEG: <BERICHT>` of `BEWEEG: <OBJECT>, (<X>, <Y>)` bevat, en niets anders!} \\
\\ 
\texttt{Jullie hebben gezamenlijk 12 van 12 straffen verzameld.} \\
\texttt{Probeer het opnieuw!} \\
            }
        }
    }
    & \\ \\

    \theutterance \stepcounter{utterance}  
    & & & \multicolumn{4}{p{0.6\linewidth}}{
        \cellcolor[rgb]{0.9,0.9,0.9}{
            \makecell[{{p{\linewidth}}}]{
                \texttt{\tiny{[GM$\langle$P2]}}
                \texttt{```} \\
\texttt{BEWEEG: C, (1,2)} \\
\texttt{```} \\
            }
        }
    }
     \\ \\

    \theutterance \stepcounter{utterance}  
    & & & \multicolumn{2}{p{0.3\linewidth}}{
        \cellcolor[rgb]{0.9,0.9,0.9}{
            \makecell[{{p{\linewidth}}}]{
                \texttt{\tiny{[GM$|$GM]}}
                \texttt{```} \\
\texttt{BEWEEG: C, (1,2)} \\
\texttt{```} \\
            }
        }
    }
    & & \\ \\

    \theutterance \stepcounter{utterance}  
    & & & \multicolumn{2}{p{0.3\linewidth}}{
        \cellcolor[rgb]{0.9,0.9,0.9}{
            \makecell[{{p{\linewidth}}}]{
                \texttt{\tiny{[GM$|$GM]}}
                \texttt{C moved to (1, 2)} \\
            }
        }
    }
    & & \\ \\

    \theutterance \stepcounter{utterance}  
    & & \multicolumn{4}{p{0.6\linewidth}}{
        \cellcolor[rgb]{0.9,0.9,0.9}{
            \makecell[{{p{\linewidth}}}]{
                \texttt{\tiny{[P1$\langle$GM]}}
                \texttt{Je bericht is doorgestuurd naar de andere speler.} \\
\\ 
\\ 
\texttt{Je speelt momenteel ronde 2 van maximaal 20.} \\
\texttt{Jullie hebben gezamenlijk 12 van 12 straffen verzameld.} \\
\texttt{De andere speler heeft een voorwerp op zijn rooster verplaatst.} \\
\\ 
\\ 
\texttt{Wat is je volgende opdracht?} \\
            }
        }
    }
    & \\ \\

    \theutterance \stepcounter{utterance}  
    & \multicolumn{4}{p{0.6\linewidth}}{
        \cellcolor[rgb]{0.9,0.9,0.9}{
            \makecell[{{p{\linewidth}}}]{
                \texttt{\tiny{[P1$\rangle$GM]}}
                \texttt{```} \\
\texttt{BEWEEG: C, (1, 2)} \\
\texttt{```} \\
            }
        }
    }
    & & \\ \\

    \theutterance \stepcounter{utterance}  
    & & & \multicolumn{2}{p{0.3\linewidth}}{
        \cellcolor[rgb]{0.9,0.9,0.9}{
            \makecell[{{p{\linewidth}}}]{
                \texttt{\tiny{[GM$|$GM]}}
                \texttt{```} \\
\texttt{BEWEEG: C, (1, 2)} \\
\texttt{```} \\
            }
        }
    }
    & & \\ \\

    \theutterance \stepcounter{utterance}  
    & & & \multicolumn{2}{p{0.3\linewidth}}{
        \cellcolor[rgb]{0.9,0.9,0.9}{
            \makecell[{{p{\linewidth}}}]{
                \texttt{\tiny{[GM$|$GM]}}
                \texttt{C moved to (1, 2)} \\
            }
        }
    }
    & & \\ \\

    \theutterance \stepcounter{utterance}  
    & & \multicolumn{4}{p{0.6\linewidth}}{
        \cellcolor[rgb]{0.9,0.9,0.9}{
            \makecell[{{p{\linewidth}}}]{
                \texttt{\tiny{[GM$\rangle$P2]}}
                \texttt{Verplaatst 'C' naar (1,2) met succes. Uw bijgewerkte raster ziet er als volgt uit:} \\
\texttt{```} \\
\texttt{ 1234567} \\
\texttt{╔═╤═╤═══╗} \\
\texttt{║H│◌│◌◌◌║ 1} \\
\texttt{║C├─┤◌◌◌║ 2} \\
\texttt{╟─┤◌├─┬─╢ 3} \\
\texttt{║◌│◌│◌│◌║ 4} \\
\texttt{╟─┼─┘◌│◌║ 5} \\
\texttt{║◌│◌◌┌┴─╢ 6} \\
\texttt{║T│WI│◌◌║ 7} \\
\texttt{╚═╧══╧══╝} \\
\texttt{```} \\
\texttt{Je speelt momenteel ronde 3 van maximaal 20.} \\
\texttt{Jullie hebben gezamenlijk 12 van 12 straffen verzameld.} \\
\texttt{De andere speler heeft een voorwerp op zijn rooster verplaatst.} \\
\\ 
\\ 
\texttt{Wat is je volgende opdracht?} \\
            }
        }
    }
    & \\ \\

    \theutterance \stepcounter{utterance}  
    & & & \multicolumn{4}{p{0.6\linewidth}}{
        \cellcolor[rgb]{0.9,0.9,0.9}{
            \makecell[{{p{\linewidth}}}]{
                \texttt{\tiny{[GM$\langle$P2]}}
                \texttt{```} \\
\texttt{BEWEEG: W, (3,2)} \\
\texttt{```} \\
            }
        }
    }
     \\ \\

    \theutterance \stepcounter{utterance}  
    & & & \multicolumn{2}{p{0.3\linewidth}}{
        \cellcolor[rgb]{0.9,0.9,0.9}{
            \makecell[{{p{\linewidth}}}]{
                \texttt{\tiny{[GM$|$GM]}}
                \texttt{```} \\
\texttt{BEWEEG: W, (3,2)} \\
\texttt{```} \\
            }
        }
    }
    & & \\ \\

    \theutterance \stepcounter{utterance}  
    & & & \multicolumn{2}{p{0.3\linewidth}}{
        \cellcolor[rgb]{0.9,0.9,0.9}{
            \makecell[{{p{\linewidth}}}]{
                \texttt{\tiny{[GM$|$GM]}}
                \texttt{Straf: (3,2) is niet leeg, maar bevat '─'.} \\
\\ 
\texttt{Jullie hebben gezamenlijk 13 van 12 straffen verzameld.} \\
\texttt{Probeer het opnieuw!} \\
            }
        }
    }
    & & \\ \\

    \theutterance \stepcounter{utterance}  
    & & \multicolumn{4}{p{0.6\linewidth}}{
        \cellcolor[rgb]{0.9,0.9,0.9}{
            \makecell[{{p{\linewidth}}}]{
                \texttt{\tiny{[GM$\rangle$P2]}}
                \texttt{Straf: (3,2) is niet leeg, maar bevat '─'.} \\
\\ 
\texttt{Jullie hebben gezamenlijk 13 van 12 straffen verzameld.} \\
\texttt{Probeer het opnieuw!} \\
            }
        }
    }
    & \\ \\

    \theutterance \stepcounter{utterance}  
    & & & \multicolumn{4}{p{0.6\linewidth}}{
        \cellcolor[rgb]{0.9,0.9,0.9}{
            \makecell[{{p{\linewidth}}}]{
                \texttt{\tiny{[GM$\langle$P2]}}
                \texttt{```} \\
\texttt{ZEG: Let op: (3,2) is een horizontale lijn op mijn bord. Kunnen we W naar (3,1) verplaatsen in plaats van (3,2)?} \\
\texttt{```} \\
            }
        }
    }
     \\ \\

    \theutterance \stepcounter{utterance}  
    & & & \multicolumn{2}{p{0.3\linewidth}}{
        \cellcolor[rgb]{0.9,0.9,0.9}{
            \makecell[{{p{\linewidth}}}]{
                \texttt{\tiny{[GM$|$GM]}}
                \texttt{```} \\
\texttt{ZEG: Let op: (3,2) is een horizontale lijn op mijn bord. Kunnen we W naar (3,1) verplaatsen in plaats van (3,2)?} \\
\texttt{```} \\
            }
        }
    }
    & & \\ \\

    \theutterance \stepcounter{utterance}  
    & & \multicolumn{4}{p{0.6\linewidth}}{
        \cellcolor[rgb]{0.9,0.9,0.9}{
            \makecell[{{p{\linewidth}}}]{
                \texttt{\tiny{[P1$\langle$GM]}}
                \texttt{Verplaatst 'C' naar (1,2) met succes. Uw bijgewerkte raster ziet er als volgt uit:} \\
\texttt{```} \\
\texttt{ 1234567} \\
\texttt{╔═╤═╤═══╗} \\
\texttt{║◌│◌│◌TI║ 1} \\
\texttt{║C├─┤W◌◌║ 2} \\
\texttt{╟─┤◌├─┬─╢ 3} \\
\texttt{║◌│◌│◌│◌║ 4} \\
\texttt{╟─┼─┘◌│◌║ 5} \\
\texttt{║◌│◌H┌┴─╢ 6} \\
\texttt{║◌│◌◌│◌◌║ 7} \\
\texttt{╚═╧══╧══╝} \\
\texttt{```} \\
\texttt{Je speelt momenteel ronde 3 van maximaal 20.} \\
\texttt{Jullie hebben gezamenlijk 13 van 12 straffen verzameld.} \\
\texttt{De andere speler stuurde dit bericht:} \\
\\ 
\\ 
\texttt{"Let op: (3,2) is een horizontale lijn op mijn bord. Kunnen we W naar (3,1) verplaatsen in plaats van (3,2)?"} \\
\\ 
\\ 
\texttt{Wat is je volgende opdracht?} \\
            }
        }
    }
    & \\ \\

    \theutterance \stepcounter{utterance}  
    & \multicolumn{4}{p{0.6\linewidth}}{
        \cellcolor[rgb]{0.9,0.9,0.9}{
            \makecell[{{p{\linewidth}}}]{
                \texttt{\tiny{[P1$\rangle$GM]}}
                \texttt{```} \\
\texttt{ZEG: Akkoord, verplaats W naar (3,1). Ik verplaats T naar (5,3) als test.} \\
\texttt{```} \\
            }
        }
    }
    & & \\ \\

    \theutterance \stepcounter{utterance}  
    & & & \multicolumn{2}{p{0.3\linewidth}}{
        \cellcolor[rgb]{0.9,0.9,0.9}{
            \makecell[{{p{\linewidth}}}]{
                \texttt{\tiny{[GM$|$GM]}}
                \texttt{```} \\
\texttt{ZEG: Akkoord, verplaats W naar (3,1). Ik verplaats T naar (5,3) als test.} \\
\texttt{```} \\
            }
        }
    }
    & & \\ \\

    \theutterance \stepcounter{utterance}  
    & & \multicolumn{4}{p{0.6\linewidth}}{
        \cellcolor[rgb]{0.9,0.9,0.9}{
            \makecell[{{p{\linewidth}}}]{
                \texttt{\tiny{[GM$\rangle$P2]}}
                \texttt{Je bericht is doorgestuurd naar de andere speler.} \\
\\ 
\\ 
\texttt{Je speelt momenteel ronde 4 van maximaal 20.} \\
\texttt{Jullie hebben gezamenlijk 13 van 12 straffen verzameld.} \\
\texttt{De andere speler stuurde dit bericht:} \\
\\ 
\\ 
\texttt{"Akkoord, verplaats W naar (3,1). Ik verplaats T naar (5,3) als test."} \\
\\ 
\\ 
\texttt{Wat is je volgende opdracht?} \\
            }
        }
    }
    & \\ \\

    \theutterance \stepcounter{utterance}  
    & & & \multicolumn{4}{p{0.6\linewidth}}{
        \cellcolor[rgb]{0.9,0.9,0.9}{
            \makecell[{{p{\linewidth}}}]{
                \texttt{\tiny{[GM$\langle$P2]}}
                \texttt{```} \\
\texttt{BEWEEG: W, (3,1)} \\
\texttt{```} \\
            }
        }
    }
     \\ \\

    \theutterance \stepcounter{utterance}  
    & & & \multicolumn{2}{p{0.3\linewidth}}{
        \cellcolor[rgb]{0.9,0.9,0.9}{
            \makecell[{{p{\linewidth}}}]{
                \texttt{\tiny{[GM$|$GM]}}
                \texttt{```} \\
\texttt{BEWEEG: W, (3,1)} \\
\texttt{```} \\
            }
        }
    }
    & & \\ \\

    \theutterance \stepcounter{utterance}  
    & & & \multicolumn{2}{p{0.3\linewidth}}{
        \cellcolor[rgb]{0.9,0.9,0.9}{
            \makecell[{{p{\linewidth}}}]{
                \texttt{\tiny{[GM$|$GM]}}
                \texttt{W moved to (3, 1)} \\
            }
        }
    }
    & & \\ \\

    \theutterance \stepcounter{utterance}  
    & & \multicolumn{4}{p{0.6\linewidth}}{
        \cellcolor[rgb]{0.9,0.9,0.9}{
            \makecell[{{p{\linewidth}}}]{
                \texttt{\tiny{[P1$\langle$GM]}}
                \texttt{Je bericht is doorgestuurd naar de andere speler.} \\
\\ 
\\ 
\texttt{Je speelt momenteel ronde 4 van maximaal 20.} \\
\texttt{Jullie hebben gezamenlijk 13 van 12 straffen verzameld.} \\
\texttt{De andere speler heeft een voorwerp op zijn rooster verplaatst.} \\
\\ 
\\ 
\texttt{Wat is je volgende opdracht?} \\
            }
        }
    }
    & \\ \\

    \theutterance \stepcounter{utterance}  
    & \multicolumn{4}{p{0.6\linewidth}}{
        \cellcolor[rgb]{0.9,0.9,0.9}{
            \makecell[{{p{\linewidth}}}]{
                \texttt{\tiny{[P1$\rangle$GM]}}
                \texttt{```} \\
\texttt{BEWEEG: T, (5, 3)} \\
\texttt{```} \\
            }
        }
    }
    & & \\ \\

    \theutterance \stepcounter{utterance}  
    & & & \multicolumn{2}{p{0.3\linewidth}}{
        \cellcolor[rgb]{0.9,0.9,0.9}{
            \makecell[{{p{\linewidth}}}]{
                \texttt{\tiny{[GM$|$GM]}}
                \texttt{```} \\
\texttt{BEWEEG: T, (5, 3)} \\
\texttt{```} \\
            }
        }
    }
    & & \\ \\

    \theutterance \stepcounter{utterance}  
    & & & \multicolumn{2}{p{0.3\linewidth}}{
        \cellcolor[rgb]{0.9,0.9,0.9}{
            \makecell[{{p{\linewidth}}}]{
                \texttt{\tiny{[GM$|$GM]}}
                \texttt{Straf: (5,3) is niet leeg, maar bevat '─'.} \\
\\ 
\texttt{Jullie hebben gezamenlijk 14 van 12 straffen verzameld.} \\
\texttt{Probeer het opnieuw!} \\
            }
        }
    }
    & & \\ \\

    \theutterance \stepcounter{utterance}  
    & & \multicolumn{4}{p{0.6\linewidth}}{
        \cellcolor[rgb]{0.9,0.9,0.9}{
            \makecell[{{p{\linewidth}}}]{
                \texttt{\tiny{[P1$\langle$GM]}}
                \texttt{Straf: (5,3) is niet leeg, maar bevat '─'.} \\
\\ 
\texttt{Jullie hebben gezamenlijk 14 van 12 straffen verzameld.} \\
\texttt{Probeer het opnieuw!} \\
            }
        }
    }
    & \\ \\

    \theutterance \stepcounter{utterance}  
    & \multicolumn{4}{p{0.6\linewidth}}{
        \cellcolor[rgb]{0.9,0.9,0.9}{
            \makecell[{{p{\linewidth}}}]{
                \texttt{\tiny{[P1$\rangle$GM]}}
                \texttt{```} \\
\texttt{ZEG: (5,3) is een horizontale lijn. Alternatieve positie voor T: (5,4). Akkoord?} \\
\texttt{```} \\
            }
        }
    }
    & & \\ \\

    \theutterance \stepcounter{utterance}  
    & & & \multicolumn{2}{p{0.3\linewidth}}{
        \cellcolor[rgb]{0.9,0.9,0.9}{
            \makecell[{{p{\linewidth}}}]{
                \texttt{\tiny{[GM$|$GM]}}
                \texttt{```} \\
\texttt{ZEG: (5,3) is een horizontale lijn. Alternatieve positie voor T: (5,4). Akkoord?} \\
\texttt{```} \\
            }
        }
    }
    & & \\ \\

    \theutterance \stepcounter{utterance}  
    & & \multicolumn{4}{p{0.6\linewidth}}{
        \cellcolor[rgb]{0.9,0.9,0.9}{
            \makecell[{{p{\linewidth}}}]{
                \texttt{\tiny{[GM$\rangle$P2]}}
                \texttt{Verplaatst 'W' naar (3,1) met succes. Uw bijgewerkte raster ziet er als volgt uit:} \\
\texttt{```} \\
\texttt{ 1234567} \\
\texttt{╔═╤═╤═══╗} \\
\texttt{║H│W│◌◌◌║ 1} \\
\texttt{║C├─┤◌◌◌║ 2} \\
\texttt{╟─┤◌├─┬─╢ 3} \\
\texttt{║◌│◌│◌│◌║ 4} \\
\texttt{╟─┼─┘◌│◌║ 5} \\
\texttt{║◌│◌◌┌┴─╢ 6} \\
\texttt{║T│◌I│◌◌║ 7} \\
\texttt{╚═╧══╧══╝} \\
\texttt{```} \\
\texttt{Je speelt momenteel ronde 5 van maximaal 20.} \\
\texttt{Jullie hebben gezamenlijk 14 van 12 straffen verzameld.} \\
\texttt{De andere speler stuurde dit bericht:} \\
\\ 
\\ 
\texttt{"(5,3) is een horizontale lijn. Alternatieve positie voor T: (5,4). Akkoord?"} \\
\\ 
\\ 
\texttt{Wat is je volgende opdracht?} \\
            }
        }
    }
    & \\ \\

    \theutterance \stepcounter{utterance}  
    & & & \multicolumn{4}{p{0.6\linewidth}}{
        \cellcolor[rgb]{0.9,0.9,0.9}{
            \makecell[{{p{\linewidth}}}]{
                \texttt{\tiny{[GM$\langle$P2]}}
                \texttt{```} \\
\texttt{ZEG: Akkoord met T(5,4). Ik verplaats nu I naar (6,3).} \\
\texttt{```} \\
            }
        }
    }
     \\ \\

    \theutterance \stepcounter{utterance}  
    & & & \multicolumn{2}{p{0.3\linewidth}}{
        \cellcolor[rgb]{0.9,0.9,0.9}{
            \makecell[{{p{\linewidth}}}]{
                \texttt{\tiny{[GM$|$GM]}}
                \texttt{```} \\
\texttt{ZEG: Akkoord met T(5,4). Ik verplaats nu I naar (6,3).} \\
\texttt{```} \\
            }
        }
    }
    & & \\ \\

    \theutterance \stepcounter{utterance}  
    & & \multicolumn{4}{p{0.6\linewidth}}{
        \cellcolor[rgb]{0.9,0.9,0.9}{
            \makecell[{{p{\linewidth}}}]{
                \texttt{\tiny{[P1$\langle$GM]}}
                \texttt{Je bericht is doorgestuurd naar de andere speler.} \\
\\ 
\\ 
\texttt{Je speelt momenteel ronde 5 van maximaal 20.} \\
\texttt{Jullie hebben gezamenlijk 14 van 12 straffen verzameld.} \\
\texttt{De andere speler stuurde dit bericht:} \\
\\ 
\\ 
\texttt{"Akkoord met T(5,4). Ik verplaats nu I naar (6,3)."} \\
\\ 
\\ 
\texttt{Wat is je volgende opdracht?} \\
            }
        }
    }
    & \\ \\

    \theutterance \stepcounter{utterance}  
    & \multicolumn{4}{p{0.6\linewidth}}{
        \cellcolor[rgb]{0.9,0.9,0.9}{
            \makecell[{{p{\linewidth}}}]{
                \texttt{\tiny{[P1$\rangle$GM]}}
                \texttt{```} \\
\texttt{BEWEEG: I, (6, 3)} \\
\texttt{```} \\
            }
        }
    }
    & & \\ \\

    \theutterance \stepcounter{utterance}  
    & & & \multicolumn{2}{p{0.3\linewidth}}{
        \cellcolor[rgb]{0.9,0.9,0.9}{
            \makecell[{{p{\linewidth}}}]{
                \texttt{\tiny{[GM$|$GM]}}
                \texttt{```} \\
\texttt{BEWEEG: I, (6, 3)} \\
\texttt{```} \\
            }
        }
    }
    & & \\ \\

    \theutterance \stepcounter{utterance}  
    & & & \multicolumn{2}{p{0.3\linewidth}}{
        \cellcolor[rgb]{0.9,0.9,0.9}{
            \makecell[{{p{\linewidth}}}]{
                \texttt{\tiny{[GM$|$GM]}}
                \texttt{Straf: (6,3) is niet leeg, maar bevat '┬'.} \\
\\ 
\texttt{Jullie hebben gezamenlijk 15 van 12 straffen verzameld.} \\
\texttt{Probeer het opnieuw!} \\
            }
        }
    }
    & & \\ \\

    \theutterance \stepcounter{utterance}  
    & & \multicolumn{4}{p{0.6\linewidth}}{
        \cellcolor[rgb]{0.9,0.9,0.9}{
            \makecell[{{p{\linewidth}}}]{
                \texttt{\tiny{[P1$\langle$GM]}}
                \texttt{Straf: (6,3) is niet leeg, maar bevat '┬'.} \\
\\ 
\texttt{Jullie hebben gezamenlijk 15 van 12 straffen verzameld.} \\
\texttt{Probeer het opnieuw!} \\
            }
        }
    }
    & \\ \\

    \theutterance \stepcounter{utterance}  
    & \multicolumn{4}{p{0.6\linewidth}}{
        \cellcolor[rgb]{0.9,0.9,0.9}{
            \makecell[{{p{\linewidth}}}]{
                \texttt{\tiny{[P1$\rangle$GM]}}
                \texttt{```} \\
\texttt{ZEG: (6,3) is een kruising. Alternatief voor I: (6,4). Akkoord?} \\
\texttt{```} \\
            }
        }
    }
    & & \\ \\

    \theutterance \stepcounter{utterance}  
    & & & \multicolumn{2}{p{0.3\linewidth}}{
        \cellcolor[rgb]{0.9,0.9,0.9}{
            \makecell[{{p{\linewidth}}}]{
                \texttt{\tiny{[GM$|$GM]}}
                \texttt{```} \\
\texttt{ZEG: (6,3) is een kruising. Alternatief voor I: (6,4). Akkoord?} \\
\texttt{```} \\
            }
        }
    }
    & & \\ \\

    \theutterance \stepcounter{utterance}  
    & & \multicolumn{4}{p{0.6\linewidth}}{
        \cellcolor[rgb]{0.9,0.9,0.9}{
            \makecell[{{p{\linewidth}}}]{
                \texttt{\tiny{[GM$\rangle$P2]}}
                \texttt{Je bericht is doorgestuurd naar de andere speler.} \\
\\ 
\\ 
\texttt{Je speelt momenteel ronde 6 van maximaal 20.} \\
\texttt{Jullie hebben gezamenlijk 15 van 12 straffen verzameld.} \\
\texttt{De andere speler stuurde dit bericht:} \\
\\ 
\\ 
\texttt{"(6,3) is een kruising. Alternatief voor I: (6,4). Akkoord?"} \\
\\ 
\\ 
\texttt{Wat is je volgende opdracht?} \\
            }
        }
    }
    & \\ \\

    \theutterance \stepcounter{utterance}  
    & & & \multicolumn{4}{p{0.6\linewidth}}{
        \cellcolor[rgb]{0.9,0.9,0.9}{
            \makecell[{{p{\linewidth}}}]{
                \texttt{\tiny{[GM$\langle$P2]}}
                \texttt{```} \\
\texttt{ZEG: Akkoord met I(6,4). Ik verplaats nu T naar (5,4).} \\
\texttt{```} \\
            }
        }
    }
     \\ \\

    \theutterance \stepcounter{utterance}  
    & & & \multicolumn{2}{p{0.3\linewidth}}{
        \cellcolor[rgb]{0.9,0.9,0.9}{
            \makecell[{{p{\linewidth}}}]{
                \texttt{\tiny{[GM$|$GM]}}
                \texttt{```} \\
\texttt{ZEG: Akkoord met I(6,4). Ik verplaats nu T naar (5,4).} \\
\texttt{```} \\
            }
        }
    }
    & & \\ \\

    \theutterance \stepcounter{utterance}  
    & & \multicolumn{4}{p{0.6\linewidth}}{
        \cellcolor[rgb]{0.9,0.9,0.9}{
            \makecell[{{p{\linewidth}}}]{
                \texttt{\tiny{[P1$\langle$GM]}}
                \texttt{Je bericht is doorgestuurd naar de andere speler.} \\
\\ 
\\ 
\texttt{Je speelt momenteel ronde 6 van maximaal 20.} \\
\texttt{Jullie hebben gezamenlijk 15 van 12 straffen verzameld.} \\
\texttt{De andere speler stuurde dit bericht:} \\
\\ 
\\ 
\texttt{"Akkoord met I(6,4). Ik verplaats nu T naar (5,4)."} \\
\\ 
\\ 
\texttt{Wat is je volgende opdracht?} \\
            }
        }
    }
    & \\ \\

    \theutterance \stepcounter{utterance}  
    & \multicolumn{4}{p{0.6\linewidth}}{
        \cellcolor[rgb]{0.9,0.9,0.9}{
            \makecell[{{p{\linewidth}}}]{
                \texttt{\tiny{[P1$\rangle$GM]}}
                \texttt{```} \\
\texttt{BEWEEG: T, (5, 4)} \\
\texttt{```} \\
            }
        }
    }
    & & \\ \\

    \theutterance \stepcounter{utterance}  
    & & & \multicolumn{2}{p{0.3\linewidth}}{
        \cellcolor[rgb]{0.9,0.9,0.9}{
            \makecell[{{p{\linewidth}}}]{
                \texttt{\tiny{[GM$|$GM]}}
                \texttt{```} \\
\texttt{BEWEEG: T, (5, 4)} \\
\texttt{```} \\
            }
        }
    }
    & & \\ \\

    \theutterance \stepcounter{utterance}  
    & & & \multicolumn{2}{p{0.3\linewidth}}{
        \cellcolor[rgb]{0.9,0.9,0.9}{
            \makecell[{{p{\linewidth}}}]{
                \texttt{\tiny{[GM$|$GM]}}
                \texttt{T moved to (5, 4)} \\
            }
        }
    }
    & & \\ \\

    \theutterance \stepcounter{utterance}  
    & & \multicolumn{4}{p{0.6\linewidth}}{
        \cellcolor[rgb]{0.9,0.9,0.9}{
            \makecell[{{p{\linewidth}}}]{
                \texttt{\tiny{[GM$\rangle$P2]}}
                \texttt{Je bericht is doorgestuurd naar de andere speler.} \\
\\ 
\\ 
\texttt{Je speelt momenteel ronde 7 van maximaal 20.} \\
\texttt{Jullie hebben gezamenlijk 15 van 12 straffen verzameld.} \\
\texttt{De andere speler heeft een voorwerp op zijn rooster verplaatst.} \\
\\ 
\\ 
\texttt{Wat is je volgende opdracht?} \\
            }
        }
    }
    & \\ \\

    \theutterance \stepcounter{utterance}  
    & & & \multicolumn{4}{p{0.6\linewidth}}{
        \cellcolor[rgb]{0.9,0.9,0.9}{
            \makecell[{{p{\linewidth}}}]{
                \texttt{\tiny{[GM$\langle$P2]}}
                \texttt{```} \\
\texttt{BEWEEG: T, (5,4)} \\
\texttt{```} \\
            }
        }
    }
     \\ \\

    \theutterance \stepcounter{utterance}  
    & & & \multicolumn{2}{p{0.3\linewidth}}{
        \cellcolor[rgb]{0.9,0.9,0.9}{
            \makecell[{{p{\linewidth}}}]{
                \texttt{\tiny{[GM$|$GM]}}
                \texttt{```} \\
\texttt{BEWEEG: T, (5,4)} \\
\texttt{```} \\
            }
        }
    }
    & & \\ \\

    \theutterance \stepcounter{utterance}  
    & & & \multicolumn{2}{p{0.3\linewidth}}{
        \cellcolor[rgb]{0.9,0.9,0.9}{
            \makecell[{{p{\linewidth}}}]{
                \texttt{\tiny{[GM$|$GM]}}
                \texttt{T moved to (5, 4)} \\
            }
        }
    }
    & & \\ \\

    \theutterance \stepcounter{utterance}  
    & & \multicolumn{4}{p{0.6\linewidth}}{
        \cellcolor[rgb]{0.9,0.9,0.9}{
            \makecell[{{p{\linewidth}}}]{
                \texttt{\tiny{[P1$\langle$GM]}}
                \texttt{Verplaatst 'T' naar (5,4) met succes. Uw bijgewerkte raster ziet er als volgt uit:} \\
\texttt{```} \\
\texttt{ 1234567} \\
\texttt{╔═╤═╤═══╗} \\
\texttt{║◌│◌│◌◌I║ 1} \\
\texttt{║C├─┤W◌◌║ 2} \\
\texttt{╟─┤◌├─┬─╢ 3} \\
\texttt{║◌│◌│T│◌║ 4} \\
\texttt{╟─┼─┘◌│◌║ 5} \\
\texttt{║◌│◌H┌┴─╢ 6} \\
\texttt{║◌│◌◌│◌◌║ 7} \\
\texttt{╚═╧══╧══╝} \\
\texttt{```} \\
\texttt{Je speelt momenteel ronde 7 van maximaal 20.} \\
\texttt{Jullie hebben gezamenlijk 15 van 12 straffen verzameld.} \\
\texttt{De andere speler heeft een voorwerp op zijn rooster verplaatst.} \\
\\ 
\\ 
\texttt{Wat is je volgende opdracht?} \\
            }
        }
    }
    & \\ \\

    \theutterance \stepcounter{utterance}  
    & \multicolumn{4}{p{0.6\linewidth}}{
        \cellcolor[rgb]{0.9,0.9,0.9}{
            \makecell[{{p{\linewidth}}}]{
                \texttt{\tiny{[P1$\rangle$GM]}}
                \texttt{```} \\
\texttt{ZEG: Mijn bord na verplaatsingen:} \\
\texttt{C(1,2), W(4,2), T(5,4), I(6,1), H(4,6).} \\
\texttt{Jouw voorstel was I(6,4), maar (6,4) is leeg bij mij.} \\
\texttt{Laten we I naar (6,4) verplaatsen en H naar (4,5)?} \\
\texttt{```} \\
            }
        }
    }
    & & \\ \\

    \theutterance \stepcounter{utterance}  
    & & & \multicolumn{2}{p{0.3\linewidth}}{
        \cellcolor[rgb]{0.9,0.9,0.9}{
            \makecell[{{p{\linewidth}}}]{
                \texttt{\tiny{[GM$|$GM]}}
                \texttt{```} \\
\texttt{ZEG: Mijn bord na verplaatsingen:} \\
\texttt{C(1,2), W(4,2), T(5,4), I(6,1), H(4,6).} \\
\texttt{Jouw voorstel was I(6,4), maar (6,4) is leeg bij mij.} \\
\texttt{Laten we I naar (6,4) verplaatsen en H naar (4,5)?} \\
\texttt{```} \\
            }
        }
    }
    & & \\ \\

    \theutterance \stepcounter{utterance}  
    & & & \multicolumn{2}{p{0.3\linewidth}}{
        \cellcolor[rgb]{0.9,0.9,0.9}{
            \makecell[{{p{\linewidth}}}]{
                \texttt{\tiny{[GM$|$GM]}}
                \texttt{Invalid format: tail is not empty: '} \\
\texttt{C(1,2), W(4,2), T(5,4), I(6,1), H(4,6).} \\
\texttt{Jouw voorstel was I(6,4), maar (6,4) is leeg bij mij.} \\
\texttt{Laten we I naar (6,4) verplaatsen en H naar (4,5)?} \\
            }
        }
    }
    & & \\ \\

    \theutterance \stepcounter{utterance}  
    & & & \multicolumn{2}{p{0.3\linewidth}}{
        \cellcolor[rgb]{0.9,0.9,0.9}{
            \makecell[{{p{\linewidth}}}]{
                \texttt{\tiny{[GM$|$GM]}}
                \texttt{Maximum number of penalties exceeded} \\
            }
        }
    }
    & & \\ \\

    \theutterance \stepcounter{utterance}  
    & & & \multicolumn{2}{p{0.3\linewidth}}{
        \cellcolor[rgb]{0.9,0.9,0.9}{
            \makecell[{{p{\linewidth}}}]{
                \texttt{\tiny{[GM$|$GM]}}
                \texttt{Initial states:} \\
\texttt{```} \\
\texttt{Player 1:} \\
\texttt{ 1234567} \\
\texttt{╔═╤═╤═══╗} \\
\texttt{║C│◌│◌TI║ 1} \\
\texttt{║◌├─┤W◌◌║ 2} \\
\texttt{╟─┤◌├─┬─╢ 3} \\
\texttt{║◌│◌│◌│◌║ 4} \\
\texttt{╟─┼─┘◌│◌║ 5} \\
\texttt{║◌│◌H┌┴─╢ 6} \\
\texttt{║◌│◌◌│◌◌║ 7} \\
\texttt{╚═╧══╧══╝} \\
\\ 
\\ 
\texttt{Player 2:} \\
\texttt{ 1234567} \\
\texttt{╔═╤═╤═══╗} \\
\texttt{║H│◌│◌◌◌║ 1} \\
\texttt{║◌├─┤◌◌◌║ 2} \\
\texttt{╟─┤◌├─┬─╢ 3} \\
\texttt{║◌│◌│C│◌║ 4} \\
\texttt{╟─┼─┘◌│◌║ 5} \\
\texttt{║◌│◌◌┌┴─╢ 6} \\
\texttt{║T│WI│◌◌║ 7} \\
\texttt{╚═╧══╧══╝} \\
\\ 
\texttt{```} \\
            }
        }
    }
    & & \\ \\

    \theutterance \stepcounter{utterance}  
    & & & \multicolumn{2}{p{0.3\linewidth}}{
        \cellcolor[rgb]{0.9,0.9,0.9}{
            \makecell[{{p{\linewidth}}}]{
                \texttt{\tiny{[GM$|$GM]}}
                \texttt{End states:} \\
\texttt{```} \\
\texttt{Player 1:} \\
\texttt{ 1234567} \\
\texttt{╔═╤═╤═══╗} \\
\texttt{║◌│◌│◌◌I║ 1} \\
\texttt{║C├─┤W◌◌║ 2} \\
\texttt{╟─┤◌├─┬─╢ 3} \\
\texttt{║◌│◌│T│◌║ 4} \\
\texttt{╟─┼─┘◌│◌║ 5} \\
\texttt{║◌│◌H┌┴─╢ 6} \\
\texttt{║◌│◌◌│◌◌║ 7} \\
\texttt{╚═╧══╧══╝} \\
\\ 
\\ 
\texttt{Player 2:} \\
\texttt{ 1234567} \\
\texttt{╔═╤═╤═══╗} \\
\texttt{║H│W│◌◌◌║ 1} \\
\texttt{║C├─┤◌◌◌║ 2} \\
\texttt{╟─┤◌├─┬─╢ 3} \\
\texttt{║◌│◌│T│◌║ 4} \\
\texttt{╟─┼─┘◌│◌║ 5} \\
\texttt{║◌│◌◌┌┴─╢ 6} \\
\texttt{║◌│◌I│◌◌║ 7} \\
\texttt{╚═╧══╧══╝} \\
\\ 
\texttt{```} \\
            }
        }
    }
    & & \\ \\

    \theutterance \stepcounter{utterance}  
    & & & \multicolumn{2}{p{0.3\linewidth}}{
        \cellcolor[rgb]{0.9,0.9,0.9}{
            \makecell[{{p{\linewidth}}}]{
                \texttt{\tiny{[GM$|$GM]}}
                \texttt{* success: False} \\
\texttt{* lose: False} \\
\texttt{* aborted: True} \\
\texttt{{-}{-}{-}{-}{-}{-}{-}} \\
\texttt{* Shifts: 2.00} \\
\texttt{* Max Shifts: 8.00} \\
\texttt{* Min Shifts: 4.00} \\
\texttt{* End Distance Sum: 14.78} \\
\texttt{* Init Distance Sum: 30.73} \\
\texttt{* Expected Distance Sum: 20.95} \\
\texttt{* Penalties: 16.00} \\
\texttt{* Max Penalties: 12.00} \\
\texttt{* Rounds: 7.00} \\
\texttt{* Max Rounds: 20.00} \\
\texttt{* Object Count: 5.00} \\
            }
        }
    }
    & & \\ \\

\end{supertabular}
}

\end{document}
